% ============================================================
%  PFC -- Práctica Fuera de Clases | Unidad 2
%  Modelado de Casos de Uso y User Stories
%  Asignatura: Ingeniería de Requisitos (ISR-401) -- UTEQ 2026-2027
%  Formato: IEEE onecolumn | Cajas de guía por integrante
%  Las líneas en MAYÚSCULAS entre corchetes son guías de
%  trabajo para cada integrante -- reemplazar con contenido real
% ============================================================
\documentclass[journal,onecolumn]{IEEEtran}

% ── Codificación ────────────────────────────────────────────
\usepackage[T1]{fontenc}
\usepackage[utf8]{inputenc}
\usepackage[spanish,es-tabla]{babel}

% ── Tablas ──────────────────────────────────────────────────
\usepackage{booktabs}
\usepackage{tabularx}
\usepackage{array}
\usepackage{multirow}
\usepackage{longtable}

% ── Figuras ─────────────────────────────────────────────────
\usepackage{graphicx}
\usepackage{float}

% ── Color negro ─────────────────────────────────────────────
\usepackage{xcolor}
\usepackage[hidelinks,colorlinks=false]{hyperref}

% ── Geometría ───────────────────────────────────────────────
\usepackage{geometry}
\geometry{a4paper, top=2.5cm, bottom=2.5cm, left=2.5cm, right=2.5cm}

% ── Encabezado ──────────────────────────────────────────────
\usepackage{fancyhdr}
\pagestyle{fancy}
\fancyhf{}
\renewcommand{\headrulewidth}{0.4pt}
\fancyhead[L]{\small PFC: MundiPets -- Unidad 2 | Ingeniería de Requisitos}
\fancyhead[R]{\small Casos de Uso y User Stories}
\fancyfoot[C]{\thepage}

% ── Espaciado ───────────────────────────────────────────────
\usepackage{parskip}
\setlength{\parskip}{4pt}
\usepackage{enumitem}
\usepackage{caption}
\usepackage{titlesec}
\captionsetup{font=small, labelfont=bf}

% ── Caja de guía (fondo gris claro, borde negro) ────────────
\usepackage{mdframed}
\mdfdefinestyle{guia}{
  linecolor=black,
  linewidth=0.6pt,
  backgroundcolor=gray!10,
  innertopmargin=6pt,
  innerbottommargin=6pt,
  innerleftmargin=8pt,
  innerrightmargin=8pt,
  skipabove=4pt,
  skipbelow=4pt
}
\renewcommand{\thesection}{\arabic{section}}
\renewcommand{\thesubsection}{\thesection.\arabic{subsection}}
\renewcommand{\thesubsubsection}{\thesubsection.\arabic{subsubsection}}

\titleformat{\section}{\normalfont\Large\bfseries}{\thesection}{1em}{}
\titleformat{\subsection}{\normalfont\large\bfseries}{\thesubsection}{1em}{}
\titleformat{\subsubsection}{\normalfont\normalsize\bfseries}{\thesubsubsection}{1em}{}
% ── Comando de caja de guía ─────────────────────────────────
% \guia{INTEGRANTE}{texto de la guía}
\newcommand{\guia}[2]{%
  \begin{mdframed}[style=guia]
    \noindent\textbf{Responsable: #1}\\[2pt]
    \textit{#2}
  \end{mdframed}
}

% ── BibTeX ──────────────────────────────────────────────────
\bibliographystyle{IEEEtran}


% ============================================================
\begin{document}
\begin{titlepage}
	\centering
	
	{\large\bfseries UNIVERSIDAD TÉCNICA ESTATAL DE QUEVEDO\\[12pt]}
	
	% Se coloca la imagen directamente sin el entorno 'figure'
	\includegraphics[width=0.3\textwidth]{uteqlogo.png}\\[12pt]
	
	{\large\bfseries FACULTAD DE CIENCIAS DE LA COMPUTACIÓN Y DISEÑO DIGITAL\\[4pt]
		CARRERA DE INGENIERÍA EN SOFTWARE}
	
	\vspace{1cm}
	
	{\normalsize\bfseries TEMA:}\\[4pt]
	{\normalsize MODELAR LOS CASOS DE USO DEL SISTEMA DEL PROYECTO FIN DE CURSO EN UNA HERRAMIENTA CASE}
	
	\vspace{1cm}
		
	{\normalsize\bfseries PROYECTO FIN DE CURSO:}\\[4pt]	
	{\normalsize MundiPets}
	
	\vspace{1cm}
	
	{\normalsize \bfseries MATERIA:}\\[4pt]
	{\normalsize INGENIERÍA DE REQUERIMIENTOS}
	
	\vspace{1cm}
	
	{\normalsize\bfseries DOCENTE:}\\[4pt]
	{\normalsize ING. GUERRERO ULLOA GLEISTON CICERÓN}
	
	\vspace{1cm}
	
	{\normalsize\bfseries INTEGRANTES:}\\[6pt]
	
	
	
	{\normalsize FUERTES ARRAES EDSON DANIEL}\\
	
	\vspace{0.3cm}
	
	{\normalsize GUTIÉRREZ ORTEGA GÉNESIS ADRIANA}\\

	\vspace{0.3cm}
	
	{\normalsize MORALES SÁNCHEZ GARY ALEJANDRO}\\
	
	\vspace{1cm}
	
	{\normalsize\bfseries CURSO:}\\[4pt]
	{\normalsize CUARTO SEMESTRE PARALELO ``B''}
	
	\vspace{1cm}
	
	{\normalsize\bfseries QUEVEDO -- LOS RÍOS -- ECUADOR}\\[4pt]
	{\normalsize\bfseries 2026 -- 2027 PPA}\\[4pt]
	{\normalsize Semana 4 -- \today}
	
\end{titlepage}


\tableofcontents
\newpage

% ============================================================
%  SECCIÓN 1 -- PORTADA Y DATOS GENERALES
% ============================================================
\section{Descripción del PFC}

\subsection{Sistema del PFC}


\newpage

% ============================================================
%  SECCIÓN 2 -- INTRODUCCIÓN
% ============================================================
\section{Introducción}

\guia{Integrante 1 -- Analista Líder}{%
Redactar una introducción de 150-200 palabras que explique:
(a)~el propósito de la práctica (modelar la funcionalidad del sistema
mediante casos de uso y user stories), (b)~la herramienta CASE seleccionada
(Enterprise Architect o StarUML) y la justificación de su elección,
(c)~el alcance cubierto en coherencia con la Entrega 1A.
Citar al menos \texttt{ISO/IEC/IEEE 29148:2018} y \texttt{Pohl (2010)}
para enmarcar la actividad dentro del proceso de IR.}

[INTRODUCCIÓN DEL INFORME]

\newpage

% ============================================================
%  SECCIÓN 3 -- ACTORES Y FRONTERA DEL SISTEMA
%  Fases 1 y 2 de la guía | Responsable: Integrante 1
% ============================================================
\section{Actores y Frontera del Sistema}

\subsection{Tabla de Actores}

\guia{Integrante 1 -- Analista Líder}{%
Derivar los actores del sistema a partir de la lista de stakeholders del PFC
(Fase~2 de la guía). Para cada actor indicar: nombre, tipo (primario /
secundario / sistema externo) y descripción de su interés u objetivo frente
al sistema. Recordar que los actores están \textbf{fuera} de la frontera;
no modelar componentes internos como actores. Incluir sistemas externos cuando
corresponda (pasarelas de pago, servicios de notificación, etc.).}

\begin{table}[H]
\caption{Actores del sistema del PFC}
\label{tab:actores}
\begin{tabularx}{\textwidth}{|c|X|l|X|}
\hline
\textbf{\#} & \textbf{Nombre del Actor} & \textbf{Tipo} & \textbf{Interés / Objetivo} \\
\hline
1 & & Primario & \\
\hline
2 & & Primario & \\
\hline
3 & & Secundario & \\
\hline
4 & & Sistema externo & \\
\hline
5 & & & \\
\hline
\end{tabularx}
\end{table}

\subsection{Frontera del Sistema}

\guia{Integrante 1 -- Analista Líder}{%
Definir la frontera del sistema en un párrafo: indicar explícitamente qué
funcionalidades quedan \textbf{dentro} del sistema y cuáles quedan
\textbf{fuera}, en coherencia con el alcance declarado en la Entrega 1A.
Esta delimitación debe ser consistente con el rectángulo de frontera
que aparece en el diagrama de la Sección~4.}

[DEFINICIÓN DE LA FRONTERA DEL SISTEMA]

\newpage

% ============================================================
%  SECCIÓN 4 -- DIAGRAMA DE CASOS DE USO
%  Fase 3 de la guía | Responsable: Integrante 2
% ============================================================
\section{Diagrama de Casos de Uso}

\guia{Integrante 2 -- Modelador UML}{%
Construir el diagrama de casos de uso en la herramienta CASE seleccionada
(Enterprise Architect o StarUML). Requisitos mínimos según la guía:
\textbf{(a)}~al menos 8 casos de uso correctamente nombrados
(verbo en infinitivo + objeto, p.ej. \textit{Registrar mascota});
\textbf{(b)}~al menos 3 actores;
\textbf{(c)}~al menos una relación \texttt{include} y una \texttt{extend}
debidamente justificadas;
\textbf{(d)}~identificadores únicos (CU-01, CU-02\ldots) registrados
como propiedad de cada elemento en la herramienta.
Asignar identificadores CU-XX a cada caso de uso.
Exportar el diagrama en alta resolución e insertarlo a continuación.}

\begin{figure}[H]
  \centering
  \fbox{\parbox{0.9\textwidth}{\centering\vspace{4cm}
    \textit{[Reemplazar con: \texttt{\textbackslash includegraphics[width=\textbackslash textwidth]\{diagrama\_cu.png\}}]}\\[6pt]
    \textit{Exportar desde la herramienta CASE en alta resolución (PNG o PDF)}
  \vspace{4cm}}}
  \caption{Diagrama de casos de uso del sistema [NOMBRE DEL PFC].}
  \label{fig:diagrama_cu}
\end{figure}

\subsection{Justificación de Relaciones UML}

\guia{Integrante 2 -- Modelador UML}{%
Para cada relación \texttt{include}, \texttt{extend} y generalización
presente en el diagrama, redactar una justificación breve (2-3 oraciones)
que explique por qué esa relación es semánticamente correcta en el dominio
del PFC. Recordar: \texttt{include} = comportamiento obligatorio reutilizado;
\texttt{extend} = comportamiento condicional u opcional (con punto de extensión);
generalización = especialización real, no forzada.}

\begin{table}[H]
\caption{Justificación de relaciones UML del diagrama}
\label{tab:relaciones}
\begin{tabularx}{\textwidth}{|l|l|X|}
\hline
\textbf{Tipo} & \textbf{Relación (CU origen $\rightarrow$ CU destino)} & \textbf{Justificación} \\
\hline
\texttt{include}  & [CU-XX] $\rightarrow$ [CU-XX] & \\
\hline
\texttt{include}  & [CU-XX] $\rightarrow$ [CU-XX] & \\
\hline
\texttt{extend}   & [CU-XX] $\rightarrow$ [CU-XX] & \\
\hline
Generalización    & [Actor/CU] $\rightarrow$ [Actor/CU] & \\
\hline
\end{tabularx}
\end{table}

\newpage

% ============================================================
%  SECCIÓN 5 -- ESPECIFICACIONES TEXTUALES DE CASOS DE USO
%  Fase 4 de la guía | Responsable: Integrante 2
% ============================================================
\section{Especificaciones Textuales de Casos de Uso}

\guia{Integrante 2 -- Modelador UML}{%
Seleccionar los tres casos de uso de mayor relevancia para el sistema
(los que concentran el valor central del PFC) y especificarlos con la
plantilla de la Fase~4 de la guía. Redactar los pasos del flujo principal
alternando actor y sistema (1. El actor solicita\ldots\ 2. El sistema
valida\ldots). No incluir decisiones de diseño ni de interfaz prematuras.
Verificar que cada especificación cumpla las características de calidad
de ISO/IEC/IEEE 29148: no ambigua, completa, singular, verificable y factible.}

% ── Caso de Uso 1 ───────────────────────────────────────────
\subsection{Especificación: [CU-XX] [Nombre del Caso de Uso]}

\begin{table}[H]
\caption{Especificación textual de [CU-XX]}
\label{tab:cu01}
\begin{tabularx}{\textwidth}{|l|X|}
\hline
\textbf{Identificador}       & CU-XX \\
\hline
\textbf{Nombre}              & [Verbo en infinitivo + objeto] \\
\hline
\textbf{Actores}             & Primario: [Actor]. Secundario: [Actor, si aplica] \\
\hline
\textbf{Precondiciones}      & \begin{enumerate}[leftmargin=*,topsep=0pt]
                                 \item [Estado requerido antes de iniciar]
                                 \item [Condición adicional, si aplica]
                               \end{enumerate} \\
\hline
\textbf{Flujo principal}     & \begin{enumerate}[leftmargin=*,topsep=0pt]
                                 \item El actor [acción].
                                 \item El sistema [respuesta].
                                 \item El actor [acción].
                                 \item El sistema [respuesta].
                                 \item [Continuar hasta completar el escenario de éxito]
                               \end{enumerate} \\
\hline
\textbf{Flujos alternativos} & \textbf{FA-1} (bifurca en paso [N]): [Descripción de variación válida]. \\
\hline
\textbf{Excepciones}         & \textbf{E-1} (bifurca en paso [N]): [Condición de error y respuesta del sistema]. \\
\hline
\textbf{Postcondiciones}     & \begin{enumerate}[leftmargin=*,topsep=0pt]
                                 \item [Estado del sistema tras ejecución exitosa]
                               \end{enumerate} \\
\hline
\textbf{Requisitos asociados} & RF-XX, RF-XX [identificadores del PFC] \\
\hline
\end{tabularx}
\end{table}

% ── Caso de Uso 2 ───────────────────────────────────────────
\subsection{Especificación: [CU-XX] [Nombre del Caso de Uso]}

\begin{table}[H]
\caption{Especificación textual de [CU-XX]}
\label{tab:cu02}
\begin{tabularx}{\textwidth}{|l|X|}
\hline
\textbf{Identificador}       & CU-XX \\
\hline
\textbf{Nombre}              & [Verbo en infinitivo + objeto] \\
\hline
\textbf{Actores}             & Primario: [Actor]. Secundario: [Actor, si aplica] \\
\hline
\textbf{Precondiciones}      & \begin{enumerate}[leftmargin=*,topsep=0pt]
                                 \item [Estado requerido antes de iniciar]
                               \end{enumerate} \\
\hline
\textbf{Flujo principal}     & \begin{enumerate}[leftmargin=*,topsep=0pt]
                                 \item El actor [acción].
                                 \item El sistema [respuesta].
                                 \item El actor [acción].
                                 \item El sistema [respuesta].
                                 \item [Continuar hasta completar el escenario de éxito]
                               \end{enumerate} \\
\hline
\textbf{Flujos alternativos} & \textbf{FA-1} (bifurca en paso [N]): [Descripción]. \\
\hline
\textbf{Excepciones}         & \textbf{E-1} (bifurca en paso [N]): [Condición de error y respuesta]. \\
\hline
\textbf{Postcondiciones}     & \begin{enumerate}[leftmargin=*,topsep=0pt]
                                 \item [Estado del sistema tras ejecución exitosa]
                               \end{enumerate} \\
\hline
\textbf{Requisitos asociados} & RF-XX, RF-XX \\
\hline
\end{tabularx}
\end{table}

% ── Caso de Uso 3 ───────────────────────────────────────────
\subsection{Especificación: [CU-XX] [Nombre del Caso de Uso]}

\begin{table}[H]
\caption{Especificación textual de [CU-XX]}
\label{tab:cu03}
\begin{tabularx}{\textwidth}{|l|X|}
\hline
\textbf{Identificador}       & CU-XX \\
\hline
\textbf{Nombre}              & [Verbo en infinitivo + objeto] \\
\hline
\textbf{Actores}             & Primario: [Actor]. Secundario: [Actor, si aplica] \\
\hline
\textbf{Precondiciones}      & \begin{enumerate}[leftmargin=*,topsep=0pt]
                                 \item [Estado requerido antes de iniciar]
                               \end{enumerate} \\
\hline
\textbf{Flujo principal}     & \begin{enumerate}[leftmargin=*,topsep=0pt]
                                 \item El actor [acción].
                                 \item El sistema [respuesta].
                                 \item El actor [acción].
                                 \item El sistema [respuesta].
                                 \item [Continuar hasta completar el escenario de éxito]
                               \end{enumerate} \\
\hline
\textbf{Flujos alternativos} & \textbf{FA-1} (bifurca en paso [N]): [Descripción]. \\
\hline
\textbf{Excepciones}         & \textbf{E-1} (bifurca en paso [N]): [Condición de error y respuesta]. \\
\hline
\textbf{Postcondiciones}     & \begin{enumerate}[leftmargin=*,topsep=0pt]
                                 \item [Estado del sistema tras ejecución exitosa]
                               \end{enumerate} \\
\hline
\textbf{Requisitos asociados} & RF-XX, RF-XX \\
\hline
\end{tabularx}
\end{table}

\newpage

% ============================================================
%  SECCIÓN 6 -- USER STORIES Y CRITERIOS DE ACEPTACIÓN
%  Fase 5 de la guía | Responsable: Integrante 3
% ============================================================
\section{User Stories y Criterios de Aceptación}

\guia{Integrante 3 -- Especialista Ágil}{%
Seleccionar tres funcionalidades del PFC asociadas a casos de uso del diagrama
(pueden ser los mismos tres de la Sección~5 u otros, dejando la trazabilidad
explícita). Para cada user story: (a)~redactarla en formato Connextra completo,
(b)~verificarla contra los criterios INVEST con análisis argumentado, y
(c)~formular al menos tres criterios de aceptación en formato Gherkin
(\textit{Dado / Cuando / Entonces}) o lista de verificación medible.
Los criterios deben cubrir: escenario de éxito, al menos un escenario
alternativo o de borde, y al menos un escenario de error.
Términos vagos sin métrica (\textit{rápido, fácil, seguro}) invalidan el criterio.
Al finalizar, registrar las user stories dentro de la herramienta CASE
(como elementos de requisito o notas vinculadas a los CU).}

% ── User Story 1 ────────────────────────────────────────────
\subsection{User Story US-01}

\textbf{Plantilla Connextra:}

\begin{mdframed}[style=guia]
\textit{US-01. Como} [rol concreto del sistema], \textit{quiero} [funcionalidad específica]
\textit{para} [beneficio o valor verificable].
\end{mdframed}

\textbf{Verificación INVEST:}

\begin{table}[H]
\caption{Verificación INVEST de US-01}
\label{tab:invest01}
\begin{tabularx}{\textwidth}{|l|l|X|}
\hline
\textbf{Criterio} & \textbf{¿Cumple?} & \textbf{Justificación} \\
\hline
Independiente  & Sí / No & \\
\hline
Negociable     & Sí / No & \\
\hline
Valiosa        & Sí / No & \\
\hline
Estimable      & Sí / No & \\
\hline
Pequeña        & Sí / No & \\
\hline
Testeable      & Sí / No & \\
\hline
\end{tabularx}
\end{table}

\textbf{Criterios de aceptación (Gherkin):}

\begin{enumerate}[leftmargin=*]
  \item \textbf{Escenario de éxito:}\\
        \textit{Dado} [contexto inicial],\\
        \textit{cuando} [acción del actor],\\
        \textit{entonces} [resultado observable y medible].

  \item \textbf{Escenario alternativo / de borde:}\\
        \textit{Dado} [contexto con condición límite],\\
        \textit{cuando} [acción del actor],\\
        \textit{entonces} [resultado esperado].

  \item \textbf{Escenario de error:}\\
        \textit{Dado} [contexto con datos inválidos o fallo],\\
        \textit{cuando} [acción del actor],\\
        \textit{entonces} [mensaje de error concreto que muestra el sistema].
\end{enumerate}

\textbf{Caso de uso asociado:} CU-XX | \textbf{Requisito(s) del PFC:} RF-XX

% ── User Story 2 ────────────────────────────────────────────
\subsection{User Story US-02}

\textbf{Plantilla Connextra:}

\begin{mdframed}[style=guia]
\textit{US-02. Como} [rol concreto del sistema], \textit{quiero} [funcionalidad específica]
\textit{para} [beneficio o valor verificable].
\end{mdframed}

\textbf{Verificación INVEST:}

\begin{table}[H]
\caption{Verificación INVEST de US-02}
\label{tab:invest02}
\begin{tabularx}{\textwidth}{|l|l|X|}
\hline
\textbf{Criterio} & \textbf{¿Cumple?} & \textbf{Justificación} \\
\hline
Independiente  & Sí / No & \\
\hline
Negociable     & Sí / No & \\
\hline
Valiosa        & Sí / No & \\
\hline
Estimable      & Sí / No & \\
\hline
Pequeña        & Sí / No & \\
\hline
Testeable      & Sí / No & \\
\hline
\end{tabularx}
\end{table}

\textbf{Criterios de aceptación (Gherkin):}

\begin{enumerate}[leftmargin=*]
  \item \textbf{Escenario de éxito:}\\
        \textit{Dado} [contexto inicial],\\
        \textit{cuando} [acción del actor],\\
        \textit{entonces} [resultado observable y medible].

  \item \textbf{Escenario alternativo / de borde:}\\
        \textit{Dado} [contexto con condición límite],\\
        \textit{cuando} [acción del actor],\\
        \textit{entonces} [resultado esperado].

  \item \textbf{Escenario de error:}\\
        \textit{Dado} [contexto con datos inválidos o fallo],\\
        \textit{cuando} [acción del actor],\\
        \textit{entonces} [mensaje de error concreto que muestra el sistema].
\end{enumerate}

\textbf{Caso de uso asociado:} CU-XX | \textbf{Requisito(s) del PFC:} RF-XX

% ── User Story 3 ────────────────────────────────────────────
\subsection{User Story US-03}

\textbf{Plantilla Connextra:}

\begin{mdframed}[style=guia]
\textit{US-03. Como} [rol concreto del sistema], \textit{quiero} [funcionalidad específica]
\textit{para} [beneficio o valor verificable].
\end{mdframed}

\textbf{Verificación INVEST:}

\begin{table}[H]
\caption{Verificación INVEST de US-03}
\label{tab:invest03}
\begin{tabularx}{\textwidth}{|l|l|X|}
\hline
\textbf{Criterio} & \textbf{¿Cumple?} & \textbf{Justificación} \\
\hline
Independiente  & Sí / No & \\
\hline
Negociable     & Sí / No & \\
\hline
Valiosa        & Sí / No & \\
\hline
Estimable      & Sí / No & \\
\hline
Pequeña        & Sí / No & \\
\hline
Testeable      & Sí / No & \\
\hline
\end{tabularx}
\end{table}

\textbf{Criterios de aceptación (Gherkin):}

\begin{enumerate}[leftmargin=*]
  \item \textbf{Escenario de éxito:}\\
        \textit{Dado} [contexto inicial],\\
        \textit{cuando} [acción del actor],\\
        \textit{entonces} [resultado observable y medible].

  \item \textbf{Escenario alternativo / de borde:}\\
        \textit{Dado} [contexto con condición límite],\\
        \textit{cuando} [acción del actor],\\
        \textit{entonces} [resultado esperado].

  \item \textbf{Escenario de error:}\\
        \textit{Dado} [contexto con datos inválidos o fallo],\\
        \textit{cuando} [acción del actor],\\
        \textit{entonces} [mensaje de error concreto que muestra el sistema].
\end{enumerate}

\textbf{Caso de uso asociado:} CU-XX | \textbf{Requisito(s) del PFC:} RF-XX

\newpage

% ============================================================
%  SECCIÓN 7 -- MATRIZ DE TRAZABILIDAD
%  Fase 6 de la guía | Responsable: Integrante 1
% ============================================================
\section{Matriz de Trazabilidad}

\guia{Integrante 1 -- Analista Líder}{%
Construir la matriz de trazabilidad requisito $\rightarrow$ caso de uso
$\rightarrow$ user story, cubriendo todos los elementos del modelo.
Verificar que: (a)~ningún caso de uso quede sin requisito de origen
(o documente el descarte con justificación); (b)~ningún requisito funcional
central quede sin caso de uso que lo realice; (c)~las tres user stories
estén vinculadas a casos de uso del diagrama.
Verificar también la coherencia terminológica entre artefactos (mismo nombre
de actores, roles y conceptos del dominio en todas las secciones).}

\begin{longtable}{|c|l|c|c|}
\caption{Matriz de trazabilidad: requisito $\rightarrow$ caso de uso $\rightarrow$ user story}
\label{tab:trazabilidad}\\
\hline
\textbf{Requisito (RF-XX)} & \textbf{Descripción breve} & \textbf{Caso de Uso} & \textbf{User Story} \\
\hline
\endfirsthead
\hline
\textbf{Requisito (RF-XX)} & \textbf{Descripción breve} & \textbf{Caso de Uso} & \textbf{User Story} \\
\hline
\endhead
RF-01 & [Descripción] & CU-XX & US-XX \\
\hline
RF-02 & [Descripción] & CU-XX & -- \\
\hline
RF-03 & [Descripción] & CU-XX & US-XX \\
\hline
RF-04 & [Descripción] & CU-XX & -- \\
\hline
RF-05 & [Descripción] & CU-XX & US-XX \\
\hline
[Agregar filas según los requisitos del PFC] & & & \\
\hline
\end{longtable}

\newpage

% ============================================================
%  SECCIÓN 8 -- CONCLUSIONES
%  Responsable: distribuidas entre los tres integrantes
% ============================================================
\section{Conclusiones}

\noindent\textit{Estructura obligatoria para cada conclusión:
(1)~afirmación clara $\rightarrow$ (2)~evidencia del trabajo realizado
$\rightarrow$ (3)~implicación para el PFC o la profesión. Mínimo 80 palabras cada una.}

\subsection*{Conclusión 1}
\addcontentsline{toc}{subsection}{Conclusión 1}

\guia{Integrante 1 -- Analista Líder}{%
Redactar una conclusión sobre la importancia de definir correctamente los
actores y la frontera del sistema. ¿Cómo influyó esta delimitación en la
calidad y completitud del diagrama de casos de uso del PFC?
Vincular con ISO/IEC/IEEE 29148:2018 o Pohl (2010).}

[CONCLUSIÓN 1]

\subsection*{Conclusión 2}
\addcontentsline{toc}{subsection}{Conclusión 2}

\guia{Integrante 2 -- Modelador UML}{%
Redactar una conclusión sobre el uso de la herramienta CASE y la semántica
de las relaciones UML (\texttt{include}/\texttt{extend}/generalización).
¿Qué errores conceptuales comunes se evitaron gracias al uso de la herramienta?
Vincular con OMG UML 2.5.1 o Pohl (2010).}

[CONCLUSIÓN 2]

\subsection*{Conclusión 3}
\addcontentsline{toc}{subsection}{Conclusión 3}

\guia{Integrante 3 -- Especialista Ágil}{%
Redactar una conclusión sobre el valor de los criterios de aceptación en
formato Gherkin para la verificabilidad de las user stories del PFC.
¿Qué dificultades se encontraron al redactar criterios medibles y cómo se resolvieron?
Vincular con Cohn (2004) o IREB CPRE FL v3.1.}

[CONCLUSIÓN 3]

\subsection*{Conclusión 4}
\addcontentsline{toc}{subsection}{Conclusión 4}

\guia{Integrante 1 -- Analista Líder}{%
Redactar una conclusión integradora sobre la trazabilidad lograda entre
requisitos, casos de uso y user stories. ¿Qué beneficios concretos aporta
esta trazabilidad al desarrollo futuro del PFC?
Citar al menos dos fuentes del deber.}

[CONCLUSIÓN 4]

\newpage

% ============================================================
%  REFERENCIAS
% ============================================================
\bibliography{referencias}

\newpage

% ============================================================
%  ANEXOS
% ============================================================
\section{Anexos}

% ── Anexo A: Evidencias de la herramienta CASE ──────────────
\subsection{Evidencias de Uso de la Herramienta CASE}
\label{app:case}

\guia{Integrante 2 -- Modelador UML}{%
Insertar capturas de pantalla del entorno de la herramienta CASE con el
proyecto del equipo abierto, mostrando: (a)~la estructura del repositorio
del modelo (paquetes y elementos); (b)~los identificadores CU-XX asignados
como propiedad de cada caso de uso; (c)~al menos una vista del modelo
en edición (no solo el diagrama exportado).
Estas capturas son evidencia obligatoria de autoría según los criterios
de admisión de la guía.}

\begin{figure}[H]
  \centering
  \fbox{\parbox{0.85\textwidth}{\centering\vspace{2.5cm}
    \textit{[Insertar captura 1: estructura del proyecto en la herramienta CASE]}
  \vspace{2.5cm}}}
  \caption{Estructura del modelo en [Enterprise Architect / StarUML].}
  \label{fig:case1}
\end{figure}

\begin{figure}[H]
  \centering
  \fbox{\parbox{0.85\textwidth}{\centering\vspace{2.5cm}
    \textit{[Insertar captura 2: propiedades de un caso de uso con su identificador CU-XX]}
  \vspace{2.5cm}}}
  \caption{Propiedad de identificador asignada al caso de uso en la herramienta.}
  \label{fig:case2}
\end{figure}

% ── Anexo B: User Stories en la herramienta CASE ────────────
\subsection{Registro de User Stories en la Herramienta CASE}
\label{app:us_case}

\guia{Integrante 3 -- Especialista Ágil}{%
Insertar captura de pantalla que evidencie el registro de las tres user
stories dentro de la herramienta CASE (como elementos de requisito o notas
vinculadas a los casos de uso correspondientes), según lo exige la Fase~5
de la guía. Esta captura es parte de las evidencias de uso del entorno
evaluadas en el criterio C6.}

\begin{figure}[H]
  \centering
  \fbox{\parbox{0.85\textwidth}{\centering\vspace{2.5cm}
    \textit{[Insertar captura: user stories registradas y vinculadas en la herramienta CASE]}
  \vspace{2.5cm}}}
  \caption{User stories US-01, US-02 y US-03 registradas en el modelo CASE.}
  \label{fig:us_case}
\end{figure}

% ── Anexo C: Declaración de Uso de IA ───────────────────────
\subsection{Declaración de Uso de Inteligencia Artificial}
\label{app:ia}

\guia{Integrante 1 -- Analista Líder}{%
Completar la tabla indicando qué herramientas de IA se usaron,
para qué tarea específica y en qué sección del informe.
El equipo es responsable de la veracidad, coherencia y originalidad
del contenido final entregado.}

\begin{table}[H]
\begin{tabularx}{\textwidth}{|l|X|l|}
\hline
\textbf{Herramienta} & \textbf{Uso específico} & \textbf{Sección(es)} \\
\hline
 &  & \\
\hline
 &  & \\
\hline
\end{tabularx}
\end{table}

\noindent El equipo es responsable de la veracidad, coherencia y
originalidad del contenido final entregado.

\end{document}
